\begin{theorem}[Transport around an oriented native plaquette]%
    \label{thm:peps_torus_plaquette_holonomy}
    \lean{TNLean.PEPS.regularWalkHolonomy_torusPlaquetteWalk,
        TNLean.PEPS.regularWalkHolonomy_torusPlaquetteWalk_reverse}
    \leanok
    \uses{thm:peps_regular_walk_holonomy}
    Let both torus periods be at least three. Assign arbitrary group
    transports to its native edges. At any plaquette, let $R_b,R_t$
    denote the rightward transports on the bottom and top edges, and let
    $U_l,U_r$ denote the upward transports on the left and right edges.
    The counterclockwise and clockwise walks based at the lower left
    corner have transports
    \begin{align}
        H_+ &= U_l^{-1}R_t^{-1}U_rR_b,\\
        H_- &= R_b^{-1}U_r^{-1}R_tU_l=H_+^{-1}.
    \end{align}
    The formulas hold at every translation, including plaquettes crossing
    either periodic seam. They are the auxiliary transport identities
    underlying flux detection and the string diagrams in
    \cite[Theorem 6.15 and \S6.6]{Schuch2010PEPS}; they assert no physical
    movement or measurement operation.
\end{theorem}

\begin{proof}\leanok
    The counterclockwise walk traverses the bottom, right, top and left
    edges in that order. The last two steps reverse the directed edge
    transports, and the first traversed edge acts first. This gives the
    first product. Reversing the entire walk gives the second product.
\end{proof}
